\documentclass[a4paper,10pt]{article}

\usepackage[utf8]{inputenc}
\usepackage[T1]{fontenc}
\usepackage[german]{babel}
\usepackage{geometry}
\usepackage{xcolor}
\usepackage{titlesec}
\usepackage{hyperref}
\usepackage{fontawesome5}
\usepackage{parskip}

% Configuration
\geometry{left=2.5cm, top=2.0cm, right=2.5cm, bottom=2.5cm}
\definecolor{primarycolor}{RGB}{0, 51, 102}
\pagestyle{empty}

\hypersetup{
    colorlinks=true,
    linkcolor=primarycolor,
    urlcolor=primarycolor,
    pdftitle={Anschreiben - Kartavya Niraj Dikshit},
}

\begin{document}

% SENDER INFO
\begin{flushright}
    \textbf{\color{primarycolor} Kartavya Niraj Dikshit} \\
    Halbmondstraße 7 \\
    91054 Erlangen \\
    \faPhone* (+49) 15510940829 \\
    \faEnvelope* \href{mailto:kartavya.dikshit@gmail.com}{kartavya.dikshit@gmail.com} \\
    \faLinkedin \href{http://www.linkedin.com/in/kartavya-dikshit}{linkedin.com/in/kartavya-dikshit} \\
    \vspace{1cm}
    19. März 2026
\end{flushright}

% RECIPIENT INFO
\begin{flushleft}
    HaCon Ingenieurgesellschaft mbH \\
    TPS Sales Logistics Team \\
    Lister Meile \\
    30161 Hannover \\
\end{flushleft}

\vspace{1cm}

% SUBJECT
\textbf{\large \color{primarycolor} Bewerbung als Werkstudent im Bereich Softwareentwicklung (Job ID: 493831)}

\vspace{0.5cm}

Sehr geehrte Damen und Herren,

mit großer Begeisterung verfolge ich die Entwicklung von Hacon als Vorreiter für digitale Mobilitätslösungen innerhalb der Siemens Mobility. Da ich täglich die Auswirkungen intelligenter Software im öffentlichen Nahverkehr erlebe, motiviert mich die Aussicht, aktiv an der Gestaltung einer effizienten und umweltfreundlichen Mobilität der Zukunft mitzuwirken. Als Masterstudent der Data Science an der Friedrich-Alexander-Universität Erlangen-Nürnberg (FAU) mit dem Schwerpunkt Künstliche Intelligenz und Machine Learning möchte ich mein technisches Know-how in Ihr TPS Sales Logistics Team einbringen.

Mein Studium an der FAU vermittelt mir die theoretischen Grundlagen, die für die Entwicklung präziser Verspätungsprognosen – einer Ihrer Kernaufgaben für diese Position – essenziell sind. In einem aktuellen Projekt habe ich bereits einen Prototyp für ein prädiktives Verspätungsmodell im Schienenverkehr entwickelt. Dabei konnte ich durch den Einsatz von Temporal Graph Convolutional Networks (T-GCN) in Python eine Vorhersagegenauigkeit von 91 \% erzielen. Diese Erfahrung ermöglicht es mir, komplexe Datenmuster im Bahnverkehr schnell zu durchdringen und in funktionale Softwarekomponenten zu übersetzen.

Neben meinen analytischen Fähigkeiten verfüge ich über fundierte Kenntnisse in der Full-Stack-Entwicklung. Durch meine Tätigkeit als freiberuflicher Entwickler und Praktikant bei Unternehmen wie Icertis bin ich bestens mit modernen Webtechnologien (React, Next.js) sowie der Backend-Entwicklung in Java und Python vertraut. Ich bin es gewohnt, in agilen Teams zu arbeiten, Web-Oberflächen für komplexe Datensätze zu optimieren und skalierbare Architekturen zu entwerfen. Die Dokumentation technischer Lösungen ist für mich dabei kein bloßer Nebenpunkt, sondern die Basis für nachhaltigen Teamerfolg.

Hacon bietet für mich die perfekte Kombination aus innovativer Start-up-Kultur und der Stabilität eines Weltkonzerns wie Siemens. Besonders die Arbeit am Train Planning System (TPS) reizt mich, da hier mathematische Präzision auf realen gesellschaftlichen Nutzen trifft. Gerne möchte ich Sie in einem persönlichen Gespräch davon überzeugen, wie ich Ihr Team bei der Entwicklung von Backend- und Frontend-Komponenten für die Mobilität von morgen unterstützen kann.

Meine Gehaltsvorstellung liegt bei einem marktüblichen Werkstudentengehalt für Informatik-Masterstudenten im Bereich von 16 bis 18 Euro pro Stunde.

Mit freundlichen Grüßen

\vspace{1cm}
Kartavya Niraj Dikshit

\end{document}